% Part: incompleteness
% Chapter: representability-in-q
% Section: introduction

\documentclass[../../../include/open-logic-section]{subfiles}

\begin{document}

\olfileid[pt]{inc}{req}{int}
\olsection{Introdução}

Os teoremas da incompletude aplicam-se a teorias nas quais fatos
básicos sobre funções computáveis podem ser expressos e provados.
Descreveremos uma teoria mínima desse tipo, chamada ``$\Th{Q}$'' (ou,
às vezes, ``$Q$ de Robinson'', em homenagem a Raphael Robinson).
Diremos o que significa uma função ser \emph{representável} em
$\Th{Q}$ e, em seguida, provaremos o seguinte:
\begin{quote}
  Uma função é representável em $\Th{Q}$ se, e somente se, é computável.
\end{quote}
Por um lado, isso nos fornece outro modelo de computabilidade. Mas
também o usaremos para mostrar que o conjunto
$\Setabs{!A}{\Th{Q} \Proves !A}$ não é decidível, reduzindo a ele o
problema da parada. Ao terminarmos, teremos provado coisas muito mais
fortes do que essa.

A linguagem de $\Th{Q}$ é a linguagem da aritmética; $\Th{Q}$ consiste
nos axiomas a seguir (a serem usados juntamente com os demais axiomas e
regras da lógica de primeira ordem com !!{identity}):
\begin{align*}
& \lforall[x][\lforall[y][(\eq[x'][y'] \lif \eq[x][y])]] \tag{$!Q_1$}\\
& \lforall[x][\eq/[\Obj 0][x']] \tag{$!Q_2$}\\
& \lforall[x][(\eq[x][\Obj 0] \lor \lexists[y][\eq[x][y']])] \tag{$!Q_3$}\\
& \lforall[x][\eq[(x + \Obj 0)][x]] \tag{$!Q_4$}\\
& \lforall[x][\lforall[y][\eq[(x + y')][(x + y)']]] \tag{$!Q_5$}\\
& \lforall[x][\eq[(x \times \Obj 0)][\Obj 0]] \tag{$!Q_6$}\\
& \lforall[x][\lforall[y][\eq[(x \times y')][((x \times y) + x)]]] \tag{$!Q_7$}\\
& \lforall[x][\lforall[y][(x < y \liff \lexists[z][\eq[(z' + x)][y]])]] \tag{$!Q_8$}
\end{align*}
Para cada número natural $n$, defina o numeral $\num{n}$ como o termo
$\Obj{0}^{\prime\prime\ldots\prime}$ no qual há, ao todo, $n$ marcas
de prima. Assim, $\num{0}$ é apenas a !!{constant}~$\Obj{0}$,
$\num{1}$ é $\Obj{0}'$, $\num{2}$ é $\Obj{0}''$ etc.

Como teoria da aritmética, $\Th{Q}$ é \emph{extremamente} fraca; por
exemplo, não se podem sequer provar fatos muito simples como
$\lforall[x][\eq/[x][x']]$ ou $\lforall[x][\lforall[y][(x + y) = (y +
    x)]]$. Mas veremos que grande parte do motivo pelo qual $\Th{Q}$ é
tão interessante está \emph{justamente} em ser tão fraca. Na verdade,
ela mal possui força suficiente para que o teorema da incompletude
valha. Outra razão de $\Th{Q}$ ser interessante é possuir um conjunto
\emph{finito} de axiomas.

Uma teoria mais forte que $\Th{Q}$ (chamada \emph{aritmética de Peano}
$\Th{PA}$) é obtida acrescentando-se a~$\Th{Q}$ um esquema de indução:
\[
(!A(\Obj 0) \land \lforall[x][(!A(x) \lif !A(x'))]) \lif \lforall[x][!A(x)]
\]
em que $!A(x)$ é uma fórmula qualquer. Se $!A(x)$ contém
!!{variable}s livres além de $x$, acrescentamos quantificadores
universais ao início para ligar todas elas (de modo que a instância
correspondente do esquema de indução seja !!a{sentence}). Por exemplo,
se $!A(x, y)$ também contém livre a !!{variable}~$y$, a instância
correspondente é
\[
\lforall[y][((!A(\Obj 0) \land \lforall[x][(!A(x) \lif !A(x'))]) \lif
  \lforall[x][!A(x)])]
\]
Usando instâncias do esquema de indução, pode-se provar muito mais a
partir dos axiomas de~$\Th{PA}$ do que dos axiomas de $\Th{Q}$. Na
verdade, é necessário bastante trabalho para encontrar afirmações
``naturais'' sobre os números naturais que não possam ser provadas na
aritmética de Peano!{}

\begin{defn}
\ollabel{defn:representable-fn}
  Dizemos que uma função $f(x_0,\ldots,x_k)$ dos números naturais nos
  números naturais é {\em representável em $\Th{Q}$} se existe uma
  fórmula $!A_f(x_0,\dots,x_k,y)$ tal que, sempre que
  $f(n_0,\dots,n_k) = m$, $\Th{Q}$ prova
\begin{enumerate}
\item\ollabel{defn:rep:a} $!A_f(\num{n_0}, \dots, \num{n_k}, \num{m})$
\item\ollabel{defn:rep:b} $\lforall[y][(!A_f(\num{n_0}, \dots,
\num{n_k}, y) \lif \num{m} = y)]$.
\end{enumerate}
\end{defn}

Há outras maneiras de enunciar a definição; por exemplo, poderíamos
exigir, de modo equivalente, que $\Th{Q}$ prove
$\lforall[y][(!A_f(\num{n_0}, \dots,
\num{n_k}, y) \liff \eq[y][\num{m}])]$.

\begin{thm}
\ollabel{thm:representable-iff-comp}
Uma função é representável em $\Th{Q}$ se, e somente se, é computável.
\end{thm}

Há duas direções na prova do teorema. A direção da esquerda para a
direita é bastante simples depois de estabelecida a aritmetização da
sintaxe. A outra direção exige mais trabalho. A ideia básica é esta:
adotamos ``recursiva geral'' como uma maneira de tornar precisa a noção
de ``computável'' e mostramos que toda função recursiva geral é
representável em~$\Th{Q}$. Lembre-se de que uma função é recursiva geral
se pode ser definida a partir de $\Zero$, da função sucessor~$\Succ$ e
das funções de projeção~$\Proj{n}{i}$ por composição, recursão primitiva
e minimização regular. Assim, uma forma de mostrar que toda função
recursiva geral é representável em~$\Th{Q}$ consiste em mostrar que as
funções básicas são representáveis e que, sempre que certas funções
são representáveis, também o são as funções delas definidas por
composição, recursão primitiva e minimização regular. Em outras
palavras, podemos mostrar que as funções básicas são representáveis e
que as funções representáveis são ``fechadas sob'' composição,
recursão primitiva e minimização regular. Isso garante que toda função
recursiva geral é representável.

O passo em que mostraríamos que as funções representáveis são fechadas
sob recursão primitiva revela-se difícil. Para evitá-lo, mostramos
primeiro que, na verdade, podemos dispensar a recursão primitiva. Isto
é, mostramos que toda função recursiva geral pode ser definida a partir
de funções básicas usando apenas composição e minimização regular. Para
isso, mostramos que a recursão primitiva pode ser realizada por uma
minimização regular específica. Contudo, para que isso funcione,
precisamos acrescentar algumas funções básicas: adição, multiplicação e
a função característica da relação de identidade~$\Char{=}$. Podemos,
então, provar o teorema mostrando que todas \emph{essas} funções
básicas são representáveis em~$\Th{Q}$ e que as funções representáveis
são fechadas sob composição e minimização regular.

\end{document}
